

\begin{frame}[fragile]{Principe général}

On définit des \alert{groupes} de lignes partageant une ou plusieurs valeurs.

\vfill

On ramène chaque groupe à une seule valeur en appliquant une \alert{fonction d'agrégation}.

\vfill

Cas le plus simple:  groupe = résultat d'une requête classique


\begin{minted}[baselinestretch=.9,
               fontsize=\small,
               framesep=2mm]{sql}
select count(*) as nombre, sum(capacité) as totalCapacité
from Logement
where type='Hôtel'
\end{minted}

\vfill

\begin{tabular}{c|c}
\textbf{nombre} & \textbf{totalCapacité}\\ \hline 
2 & 168\\ \hline 
\end{tabular}

\end{frame}

%%%%%%ù
\begin{frame}[fragile]{Le \texttt{group by}}

La clause \sqlkw{group by} \texttt{att1, ..., attn} \alert{partitionne} le résultat d'un
\sqlkw{select from where} en fonction des \texttt{att1, ..., attn}

\vfill 

Chaque \alert{groupe} contient les lignes qui partagent les mêmes valeurs pour \texttt{att1, ..., attn}.

\vfill

\begin{minted}[baselinestretch=.9,
               fontsize=\small,
               framesep=2mm]{sql}
select type, count(*) as Nombre, 
           sum(capacité)as totalCapacité
from Logement
group by type
\end{minted}

\vfill
Procède en deux étapes: d'abord on groupe, puis on agrège.
\end{frame}


%%%%%%ù
\begin{frame}[fragile]{Décomposons\,: l'étape de regroupement}

On obtient une structure intermédiaire, avec autant de lignes que de valeurs
distinctes pour les attributs de regroupement (ici, \texttt{type}).

\vfill
\begin{footnotesize}
\begin{tabular}{c|c}
\textbf{type} & Groupe (Logement) \\ \hline 
Auberge & \begin{tabular}{c} 
            (ca, Causses,    45,  Auberge,     Cévennes) \\ 
       \end{tabular}                
       \\ \hline 
Gîte &  \begin{tabular}{c} 
            (pi,     U Pinzutu,   10,  Gîte,    Corse)\\ 
       \end{tabular}                \\ \hline 
Hôtel &  \begin{tabular}{c} 
            (ge,     Génépi,  134,     Hôtel,   Alpes)\\
             (ta,     Tabriz,  34,  Hôtel,   Bretagne) \\ 
       \end{tabular}                \\ \hline 
\end{tabular}
\end{footnotesize}
\vfill

\alert{Ce n'est pas une table en première forme normale (atomicité
des valeurs).}
\end{frame}

%%%%%%ù
\begin{frame}[fragile]{L'étape d'agrégation}

La fonction d'agrégation ramène chaque groupe à une valeur
\vfill

\begin{tabular}{c|c|c}
\textbf{type} & \textbf{count(*)} & \textbf{sum(capacité)}\\ \hline 
Auberge & 1 & 45\\ \hline 
Gîte & 1 & 10\\ \hline 
Hôtel & 2 & 168\\ \hline 
\end{tabular}

\vfill
\alert{Cette fois c'est une table en première forme normale.}
\end{frame}


%%%%%%ù
\begin{frame}[fragile]{La clause \sqlkw{having}}

Filtrage dans un deuxième temps, après calcul des agrégats, 
portant sur le résultat de la fonction d'agrégation.
\vfill

\alert{Bien distinguer de la clause \sqlkw{where} qui s'applique aux nuplets}

\vfill

\begin{minted}[baselinestretch=.9,
               fontsize=\small,
               framesep=2mm]{sql}
select type, count(*) as nombre, 
           sum(capacité) as totalCapacité
from Logement
group by type
having count(*) > 1
\end{minted}

\vfill

\begin{tabular}{c|c|c}
\textbf{type} & \textbf{nombre} & \textbf{totalCapacité}\\ \hline 
Hôtel & 2 & 168\\ \hline 
\end{tabular}

\end{frame}

%%%%%%%%%
\begin{frame}{À retenir}

Agrégats = extension de SQL

\vfill

\begin{redbullet}
  \item  S'applique au \alert{résultat} d'une requête standard
  \item \alert{Partitionne} en groupes de nuplets partageant les mêmes valeurs de regroupement
  \item \alert{Réduit} chaque groupe à une valeur atomique grâce à une fonction d'agrégation
  \item On peut \alert{filtrer} les groupes obtenus avec \sqlkw{having}
\end{redbullet}


\end{frame}
